\heading{数学物理方法（下）-第7次作业}


\begin{question}
设有半径为 \(a\) 的导体半球，球面温度为 \(u_0\sin^2\theta\)，底面绝热，求半球内稳态温度分布。
\begin{solution}
轴对称球坐标中
\[
\nabla^2u=0,\qquad 0<r<a,\quad 0<\theta<\frac\pi2.
\]
底面绝热为 \(u_\theta(r,\pi/2)=0\)。因此只取偶阶 Legendre 多项式：
\[
u(r,\theta)=\sum_{m=0}^{\infty}A_{2m}r^{2m}P_{2m}(\cos\theta).
\]
边界 \(\sin^2\theta=1-\cos^2\theta\) 可写成
\[
\sin^2\theta=\frac23\left(P_0(\cos\theta)-P_2(\cos\theta)\right).
\]
故
\[
u(r,\theta)=\frac{2u_0}{3}
\left[1-\left(\frac ra\right)^2P_2(\cos\theta)\right].
\]
\end{solution}
\end{question}

\begin{question}
将若干函数按 Legendre 多项式展开，并计算若干含 Legendre 多项式的积分。
\begin{solution}
Legendre 展开的统一公式为
\[
f(x)=\sum_{l=0}^{\infty}a_lP_l(x),\qquad
a_l=\frac{2l+1}{2}\int_{-1}^{1}f(x)P_l(x)\,dx.
\]
这个公式来自正交性。把展开式两端同乘 \(P_k(x)\) 后积分，除 \(l=k\) 外各项为零，故
\[
\int_{-1}^{1}f(x)P_k(x)\,dx
=a_k\int_{-1}^{1}P_k^2(x)\,dx
=a_k\frac{2}{2k+1}.
\]
若函数只定义在 \(0\le x\le1\)，可先按题意作奇延拓、偶延拓或零延拓，再用同一公式。常用结果有
\[
x^3=\frac25P_1(x)+\frac35P_3(x),
\]
这是因为
\[
P_1(x)=x,\qquad P_3(x)=\frac12(5x^3-3x),
\]
解出 \(x^3=(2P_3+3P_1)/5\) 即得。又如奇函数 \(\operatorname{sgn}x\) 只含奇阶项，
\[
\operatorname{sgn}x=\sum_{\substack{l=1\\ l\ {\rm odd}}}^{\infty}
\frac{2l+1}{2}\left(\int_0^1P_l(x)\,dx-\int_{-1}^0P_l(x)\,dx\right)P_l(x).
\]
积分题主要使用正交性
\[
\int_{-1}^{1}P_k(x)P_l(x)\,dx=\frac{2}{2l+1}\delta_{kl},
\]
以及 Legendre 方程的 Lagrange 恒等式。特别地，
\[
\int_{-1}^{1}P_l(x)\ln(1-x)\,dx
=
\begin{cases}
2\ln2-2,&l=0,\\
-\dfrac{2}{l(l+1)},&l\ge1.
\end{cases}
\]
\end{solution}
\end{question}

\begin{question}
设函数 \(F(x,t)=G(2xt-t^2)\) 可在 \(t=0\) 展开为
\[
G(2xt-t^2)=\sum_{n=0}^{\infty}g_n(x)t^n.
\]
证明 \(g_n(-x)=(-1)^ng_n(x)\)，并证明递推关系
\[
g_0'(x)=0,\qquad xg_n'(x)-ng_n(x)=g'_{n-1}(x).
\]
\begin{solution}
由
\[
G(-2xt-t^2)=G(2x(-t)-(-t)^2)
\]
比较 \(t^n\) 系数即得
\[
g_n(-x)=(-1)^ng_n(x).
\]
对展开式关于 \(x,t\) 求导：
\[
F_x=2tG'(2xt-t^2),\qquad
F_t=(2x-2t)G'(2xt-t^2).
\]
消去 \(G'\) 得
\[
(x-t)F_x=tF_t.
\]
代入级数并比较 \(t^n\) 系数：
\[
xg_n'-g_{n-1}'=ng_n,
\]
即所需递推关系。于是 \(g_n\) 为不高于 \(n\) 次的多项式。
\end{solution}
\end{question}

\begin{question}
有一半径为 \(R\) 的接地理想导体球壳，外部放有半径为 \(a\)、总电量为 \(Q\) 的均匀带电圆环，\(a>R\)，环心与球心重合。求球外的电势分布。
\begin{solution}
圆环位于赤道平面时，其在球心附近的势可展开为
\[
\Phi_{\rm ext}(r,\theta)=\frac{Q}{4\pi\varepsilon_0}
\sum_{l=0}^{\infty}\frac{r^l}{a^{l+1}}P_l(0)P_l(\cos\theta),
\qquad r<a.
\]
接地球面要求在 \(r=R\) 上总势为零。像解为
\[
\Phi_{\rm im}(r,\theta)=
-\frac{Q}{4\pi\varepsilon_0}
\sum_{l=0}^{\infty}\frac{R^{2l+1}}{a^{l+1}r^{l+1}}
P_l(0)P_l(\cos\theta).
\]
故球外且环内区域的电势为两者之和；环外区域再加上相应外部多极展开即可。该式满足 \(r=R\) 接地条件。
\end{solution}
\end{question}

\begin{question}
半径为 \(a\) 的接地理想导体球壳外有一单位点电荷，点电荷与球心相距 \(d>a\)，求球外电势分布。
\begin{solution}
由电像法，像电荷位于同一直线上，距离球心
\[
d'=\frac{a^2}{d},
\]
电量为
\[
q'=-\frac ad.
\]
若真实电荷位于 \(\bm r_0\)，则球外电势为
\[
\Phi(\bm r)=\frac1{4\pi\varepsilon_0}
\left(
\frac1{|\bm r-\bm r_0|}
-\frac ad\frac1{|\bm r-\bm r_0'|}
\right),
\qquad |\bm r|>a,
\]
其中 \(\bm r_0'=(a^2/d^2)\bm r_0\)。在 \(r=a\) 上两项正好相消。
\end{solution}
\end{question}

